\begin{lemma}[The integral specialization of the two-period cusp]
\label{iv:rank_two_specialization}
\StatementRecord{iv:rank_two_specialization}{proved}{No earlier separately identified check is recorded. Checked only in this writing. Pre-continuation writing events: GPT-6 Astra (OpenAI), ChatGPT Pro mode, effort Pro, S6\_10\_PAPERS, 18 November 2026 04:46:34.548185+00:00: computation or formal certificate; targeted proof check. GPT-6 Astra (OpenAI), ChatGPT Pro mode, effort Pro, S6\_10\_PAPERS, 18 November 2026 04:50:08.441928+00:00: targeted proof check. Continuation metadata events: GPT-6 Astra (OpenAI), ChatGPT Pro mode, effort Pro, S6\_10\_PAPERS continuation, 19 November 2026: used and rechecked in this run; status-vocabulary reconciliation; metadata reconciliation only; no new mathematical check. \emph{Pre-continuation recorded proof scope.} Checked in this writing by GPT-6 Astra (OpenAI), ChatGPT Pro mode, effort Pro, 18 November 2026}{The explicit toric construction below; \cite[Lemma 7.4, pp.~18--19, cusp clause]{Engel2026}; \cite[Propositions 2.7--2.8, pp.~4--6]{Wang2026}}
For each unit-square, one-component rank-two Mumford filling used in
\cref{night:P3F.20,night:P3F.21}, the integral specialization
isomorphism has the explicit generators and unit-preserving family
map constructed below.  Its cross generator is the Chern class of
the actual Cartier line with period coefficient
\(z_1^{\nu_2}z_2^{\nu_1}q^{\nu_1\nu_2}\).
The existence of the integral isomorphism is the cited local result;
this proof supplies the explicit form required by the present
attachments.  Writing its central fibre as \(W\), its cohomology is
free with ranks
\[
                 (b_0,b_1,b_2,b_3,b_4)=(1,2,4,2,1).
\]
The assertion retains the holomorphic units in both actual period maps.
\end{lemma}
\begin{proof}
\begin{enumerate}[label=\arabic*.]
\item On the actual punctured local model the periods have the form
\[
 H(z_1,z_2,q)=(q u_1(q)z_1,u_2(q)z_2,q),\qquad
 G(z_1,z_2,q)=(u_3(q)z_1,q u_4(q)z_2,q),
\]
where the four \(u_i\) are holomorphic units on a disc.  Choose their
holomorphic logarithms \(L_i\) there and replace \(u_i\) by
\(\exp(tL_i)\), \(0\leq t\leq1\).  This defines an actual family
of the same toric period quotients, including their central fibres.
The period valuations and fan are unchanged.  On a sufficiently
small closed \(q\)-disc, the \(L_i\) and all their \(t\)-derivatives
are bounded.  The two \(\log|q|\) diagonal terms therefore dominate
the bounded matrix of unit terms uniformly in \(t\).  The usual
closed period parallelogram is a bounded fundamental region in
logarithmic coordinates.  Adding its finitely many toric faces gives
a compact fundamental region over the closed disc.  The local
equations remain \(q=x_0x_1x_2\), and all changes of chart depend
smoothly on \(t\).

\item This unit deformation gives the required morphism of
specialization diagrams, as follows.  In a toric chart take the lift
of \(\partial_t\) that holds its three monomial coordinates fixed.
It preserves \(q\) and every branch stratum.  Patch these local lifts
with a smooth partition of unity on the actual quotient.  Their
projections are all \(\partial_t\), their \(q\)-derivatives are zero,
and each is tangent to every branch through its domain; these three
properties survive the patching.  Compactness of the preceding
fundamental regions makes the resulting vector field complete for
\(0\leq t\leq1\).  Its flow is an invertible, \(q\)-preserving,
stratum-preserving map from the unit-free family to the actual
family.  The inverse is the backwards flow.  It identifies the
central cohomology, the nearby cohomology, their monodromies, and
the specialization maps, since it is a map of the entire families
over the disc.  Thus the following integral computation transports
through an actual family map.

\item We use the antidiagonal unit-square triangulation.  For the
other diagonal the explicit character map
\((z_1,z_2,q)\mapsto(z_1,z_2^{-1},q)\), together with replacement
of the second period generator by its inverse, carries its fan to
this one.  This monomial map extends on the corresponding toric
charts and its inverse is the same formula.  In the unit-free model
the normalization of \(W\) is the smooth toric surface with cyclic
rays
\[
 r_0=(1,0),\ r_1=(0,1),\ r_2=(-1,1),\quad
 r_3=-r_0,\ r_4=-r_1,\ r_5=-r_2.
\]
Consecutive rays have determinant one and
\(r_{i-1}+r_{i+1}=r_i\).  Consequently its six boundary curves
\(D_i\) have square \(-1\).  Opposite curves are identified by
the actual period translations.  There are three normalized double
curves and two normalized triple points.  The dual complex is the
one-vertex, three-edge, two-triangle triangulation of \(T^2\).

\item Here are all the integral differentials in the normalization
complex.  In degree zero its terms are
\(\mathbb Z\to\mathbb Z^3\to\mathbb Z^2\), with matrices
\[
 0,\qquad
 C_0=\begin{pmatrix}1&-1&1\\-1&1&-1\end{pmatrix}.
\]
The rows are the boundaries of the two oriented triangles.
The Picard group of the normalization is \(\mathbb Z^6\) modulo
the two character relations.  Subtracting the unique character
with prescribed values on \(r_0,r_1\) sets \(a_0=a_1=0\);
\((a_2,a_3,a_4,a_5)\) is an integral basis.  The degree on
\(D_i\) is \(a_{i-1}+a_{i+1}-a_i\).  Hence opposite-curve
restriction in degree two is exactly
\[
 C_2=\begin{pmatrix}
 -1&1&-1&1\\
  1&-1&1&-1\\
 -1&1&-1&1
 \end{pmatrix}:\mathbb Z^4\longrightarrow\mathbb Z^3.
\]
In degree four the only term is \(\mathbb Z\) from the
normalization.  Both displayed nonzero matrices have rank one
and a unit Smith invariant.  Thus their kernels and cokernels
are free.  Local exactness of the normalization complex is the
augmented simplex on the branches of \(x_0x_1x_2=0\); this also
specifies its maps at the two triple points.

\item The nonzero normalization spectral-sequence terms are in
columns \(0,1,2\) and even rows \(0,2,4\).  A differential of
length two has odd-row target, and one of length at least three
has column beyond two.  Every higher differential is therefore
zero.  The free terms in each total degree have no torsion
extension, giving exactly \((1,2,4,2,1)\).  We next determine the
integral image of specialization; these ranks alone would not
determine its index.

\item Choose a real positive nearby \(q\), write
\(\tau=\log q/(2\pi i)\), and parametrize the unit-free torus by
\[
 z_i=\exp(2\pi i(t_i+\tau s_i)),\qquad
 (s_i,t_i)\in\mathbb R^2/\mathbb Z^2.
\]
Put \(a_i=ds_i\) and \(b_i=dt_i\) as integral cohomology
classes.  A positive \(q\)-loop changes the marked period by one
angular cycle; its pullback action is
\[
                   T(a_i)=a_i,\qquad T(b_i)=b_i+a_i.
\]
Direct expansion in the ordered exterior basis gives the following
integral bases of the invariant lattices:
\[
\begin{array}{c|l}
0&1\\
1&a_1,a_2\\
2&a_1a_2,\ a_1b_1,\ a_2b_2,\ a_1b_2+a_2b_1\\
3&a_1a_2b_1,\ a_1a_2b_2\\
4&a_1b_1a_2b_2.
\end{array}
\]
For example the two nontrivial degree-two equations require the
coefficient of \(b_1b_2\) to vanish and the two cross coefficients
to be equal.  Their coefficients are units, so this is an integral
basis, with no saturation factor.

\item The identity map of the valuation lattice refines the square
fan to the stated triangulation.  The corresponding toric maps,
with their unchanged period characters, give an actual morphism
from this family to the fibre product of the two nodal elliptic
families.  It is the identity on each nearby torus.  On the central
fibre it gives \(W\to C\times C\), where
\(C\) is the nodal rational curve.  The loop around its normalized
ends gives \(H^1(C;\mathbb Z)\); the pullback of that generator
specializes to \(a_i\).  On each factor choose a smooth point of
\(C\).  Smoothness at that point gives a section on a smaller
disc by a local submersion coordinate; its image is a Cartier
divisor on the total nodal family.  The inverse divisor line
restricts to \(\mathcal O_C(-p)\).  Its specialization has
degree \(-1\) in the complex orientation.  Since the complex
orientation of \(t_i+\tau s_i\) is \(dt_i\wedge ds_i\), its
Chern class in the present marking is exactly \(a_ib_i\).
The actual pullbacks of these two lines therefore give the second
and third degree-two generators.  Cup products already give
\(a_1a_2\), both degree-three generators up to the displayed
exterior signs, and the degree-four generator with coefficient one.

\item We construct the remaining cross class with its full
equivariance.  On the lattice \(\mathbb Z^2\) put
\[
       \phi(v)=v_1v_2,\qquad
       B=\begin{pmatrix}0&1\\1&0\end{pmatrix}.
\]
Interpolate \(\phi\) affinely on every unit triangle.  Each
triangle is unimodular, so this interpolation has integral
linear coefficients on its height-one cone.  The ray coefficients
\(-\phi(v)\) define a locally finite toric Cartier divisor on
the infinite model, and thus an invertible line there.  For every
period vector \(\nu\in\mathbb Z^2\), the exact identity is
\[
 \phi(v+\nu)-\phi(v)=v^{\mathsf T}B\nu+\nu_1\nu_2.
\]
The line lift of that period is
\[
 (z_1,z_2,q,w)\longmapsto
 (q^{\nu_1}z_1,q^{\nu_2}z_2,q,
      z_1^{\nu_2}z_2^{\nu_1}q^{\nu_1\nu_2}w).
\]
On each toric chart the preceding height identity identifies the
source Cartier frame with the target Cartier frame; in those
frames its coefficient is one.  It therefore extends and is
invertible across the central boundary.  The complete cocycle
identity is
\[
 m_\nu(q^\mu z,q)m_\mu(z,q)=m_{\nu+\mu}(z,q),\qquad
 m_\nu=z^{B\nu}q^{\nu_1\nu_2},
\]
because
\(\mu^{\mathsf T}B\nu+\phi(\mu)+\phi(\nu)
=\phi(\mu+\nu)\).  Hence this is an actual descended line on
the whole toric quotient, including \(W\).

\item Its Chern class on the nearby torus is explicit.  The norm
\[
 |w|\exp\!\bigl(\pi\operatorname{Im}(\tau)
                      s^{\mathsf T}Bs\bigr)
\]
is invariant under the full line-period action.  The resulting
unitary coefficient at \(s\mapsto s+\nu\) is
\(\exp(2\pi i\nu^{\mathsf T}Bt)\).  In this unitary frame the
connection
\[
 \mathcal A=-2\pi i\,s^{\mathsf T}B\,dt
\]
has exactly the required transformation
\(\mathcal A(s+\nu)=\mathcal A(s)-m_\nu^{-1}dm_\nu\).
With \(c_1=(i/2\pi)F_{\mathcal A}\) its first Chern class is
\[
                a_1b_2+a_2b_1.
\]
This gives the last degree-two generator with coefficient one.
Every displayed invariant generator is therefore the specialization
of an actual central cohomology class.  The source and target are
free of the same ranks already computed, so this surjection is an
integral isomorphism in every degree.  Transport by the complete
unit-deformation flow in steps 1--2 proves the assertion for the
two original filled families, with their actual unit factors.
\end{enumerate}
\end{proof}
